\section{\label{sec:concl}Conclusions and future work}

We dressed up rank polymorphism in dependent types, thereby obtaining a
haute-couture collection of tools for constructing and formally reasoning about
array-based algorithms.  This collection comprises: an \emph{array theory} that
enables the definition of specifications without committing to any particular
array implementation; the \emph{initial model} of the array theory that supplies
recursion over array shapes; \emph{array combinators} defined as interfaces
with fixed semantics, ensuring that efficient implementations can be
instantiated for all models; and \emph{inductive reshapes} that make it
possible to prove theorems concerning the stability of functions
under reshapes.

We have demonstrated that our specification of \texttt{scan} can be transformed
into algorithms exhibiting differing runtime characteristics merely by
selecting a shape for the input array.  The main theorem guarantees that
applying \texttt{scan} to any reshaping of the input array that is admissible under
\AD{Reshape} yields identical values.  We have shown how monads can be employed
to model runtime properties such as sharing and parallelism of array
expressions.  Coupled with the freedom to choose a runtime representation of
arrays and to define efficient combinators for that representation, this
approach mitigates guesswork when selecting the optimal shape for a given
hardware architecture.  An additional benefit of the dependently-typed
framework is that all computations are index-safe and all functions are
provably terminating.

By no means do we claim that the given array framework is complete for expressing
every conceivable numerical computation.  Nevertheless, this work offers a
perspective on handling collections that can be represented as containers
possessing non-trivial shape structure.

The notion of arrays can be generalised by extending shapes with new
constructors such as addition, exponentiation, or perhaps even dependent sums
or products. By doing so we obtain greater flexibility in grouping elements
within arrays, although rank-polymorphic functions and reshapes become
correspondingly more intricate. The language of reshapes is not immutable---it
may be extended (\eg{} by permitting permutations) or restricted (\eg{} by
forbidding flattening) according to the problem under consideration. The array
combinators presented here remain universal provided that shapes include
singleton and product constructors. Moreover, the theory can be readily
expanded with additional array combinators without jeopardising the existing
development.

The present contribution constitutes a first step toward a wealth of promising
future work.  Firstly, we aim to generate high-performance code from the
proposed specifications.  While it is unlikely that Agda alone can produce code
competitive with highly optimised low-level libraries, it is straightforward to
encapsulate the array definitions within a small DSL and perform code
generation, as we have done for convolutional neural
networks~\citep{10.1145/3747524}.  Another substantial question concerns the
translation of scan-like algorithms into in-place array updates as described
in Blelloch's original
paper~\citep{GEBlelloch:prefixsums}.  On a theoretical level, it would be
interesting to investigate whether incorporating reshape equalities into shape
definitions simplifies specifications; however, this would necessitate moving
to more expressive type theories that support quotient types.
